% ---------------------------------------------------------------------------
% Geometria da direção de Ackermann.
%   - distância entre eixos L, bitola d;
%   - rodas traseiras fixas, rodas dianteiras esterçadas (interna/externa);
%   - centro instantâneo de rotação (ICR) e raio R do eixo traseiro;
%   - ângulo phi do modelo bicicleta (roda dianteira virtual central).
% ---------------------------------------------------------------------------
\begin{tikzpicture}[>=stealth, scale=1.1]
  % ICR
  \coordinate (O) at (0,0);
  \fill (O) circle (1.5pt) node[left]{ICR};

  % eixo traseiro (a uma distância R do ICR)
  \def\R{3.2}
  \def\L{2.2}
  \def\d{1.4}
  \coordinate (RL) at (\R,0);            % roda traseira esquerda (interna)
  \coordinate (RR) at (\R+\d,0);         % roda traseira direita (externa)
  \coordinate (RC) at (\R+0.5*\d,0);     % centro eixo traseiro
  \coordinate (FC) at (\R+0.5*\d,\L);    % centro eixo dianteiro

  % chassi
  \draw[thick, fill=blue!6] (RL) -- (RR) -- ($(RR)+(0,\L)$) -- ($(RL)+(0,\L)$) -- cycle;

  % rodas traseiras (fixas)
  \draw[line width=4pt] ($(RL)+(0,-0.25)$) -- ($(RL)+(0,0.25)$);
  \draw[line width=4pt] ($(RR)+(0,-0.25)$) -- ($(RR)+(0,0.25)$);

  % raios do ICR até cada roda dianteira
  \draw[dashed, gray] (O) -- ($(RL)+(0,\L)$);
  \draw[dashed, gray] (O) -- ($(RR)+(0,\L)$);
  \draw[dashed] (O) -- (RC);
  \draw[dashed] (O) -- (FC);

  % rodas dianteiras esterçadas (perpendiculares ao raio)
  \draw[line width=4pt, red]  ($(RL)+(0,\L)$) ++(-28:-0.25) -- ++(-28:0.5);  % interna
  \draw[line width=4pt, red]  ($(RR)+(0,\L)$) ++(-20:-0.25) -- ++(-20:0.5);  % externa

  % rótulos de dimensões
  \draw[<->] ($(RL)+(-0.4,0)$) -- ($(RL)+(-0.4,\L)$) node[midway,left]{$L$};
  \draw[<->] (RL) ++(0,-0.6) -- ($(RR)+(0,-0.6)$) node[midway,below]{$d$};
  \draw[<->] (O) ++(0,0.15) -- ($(RC)+(0,0.15)$) node[midway,above]{$R$};

  % angulos
  \node[red] at ($(RL)+(0.55,\L+0.35)$) {$\phi_i$};
  \node[red] at ($(RR)+(0.55,\L+0.35)$) {$\phi_o$};
\end{tikzpicture}
